\documentclass[10pt]{article}
\usepackage[letterpaper,margin=0.75in]{geometry}
\usepackage{amsmath,amssymb,amsfonts,mathtools}
\usepackage{titlesec}
\usepackage{enumitem}
\usepackage{graphicx}
\usepackage[hidelinks]{hyperref}

\titleformat{\section}{\large\bfseries}{\thesection.}{0.5em}{}
\titleformat{\subsection}{\normalsize\bfseries}{\thesubsection.}{0.5em}{}

\begin{document}

\title{\textbf{Logarithmic Error-Bounded Latency Aggregation: \\ Mathematical Proof, DDSketch Convergence, and Memory Complexity}}
\author{Inference Systems Engineering \& Quality Assurance Team}
\date{June 17, 2026}
\maketitle

\begin{abstract}
This report provides a formal mathematical analysis of the logarithmic binning DDSketch strategy implemented in the LLM Observability Platform for Latency and SLO tracking (\textbf{F-L-01} through \textbf{F-L-06}). We define the logarithmic mapper mapping real-valued streaming latencies to compression indices, prove the relative error bounds, define the Tokens-Per-Second (TPOT) formulation, and demonstrate that the memory complexity of the sketch is $O(1)$ with respect to request volume.
\end{abstract}

\section{Introduction}
Real-time tracking of Large Language Model (LLM) latencies is crucial for maintaining service quality. The platform consumes raw span events and aggregates them into percentile representations. Since keeping raw arrays of latencies consumes infinite memory over long windows, we employ the DataDog DDSketch algorithm to approximate percentiles with a guaranteed relative error limit $\epsilon$.

---

\section{Mathematical Model of DDSketch Logarithmic Binning}

\subsection{Logarithmic Mapping Formulation}
Let $v > 0$ be a real-valued latency sample (in milliseconds). The mapping function $I(v)$ converts $v$ into a bucket index:
\begin{equation}
I(v) = \lfloor \log_{\gamma} (v) \rfloor
\end{equation}
where the base parameter $\gamma > 1$ represents the growth factor of the bucket boundaries and is directly related to the maximum relative accuracy error $\epsilon$:
\begin{equation}
\gamma = \frac{1 + \epsilon}{1 - \epsilon}
\end{equation}

For a target relative accuracy of $\epsilon = 0.01$ (1\% error limit), the base $\gamma$ is:
\begin{equation}
\gamma = \frac{1.01}{0.99} \approx 1.020202
\end{equation}

\subsection{Proof of Relative Error Bounds}
Let $v$ be the actual value mapped to index $i = I(v)$. The boundaries of bucket $i$ are defined as:
\begin{equation}
\text{lower}(i) = \gamma^i, \quad \text{upper}(i) = \gamma^{i+1}
\end{equation}
The sketch approximates any value mapped to bucket $i$ using the bucket center or lower boundary representation:
\begin{equation}
\hat{v} = \gamma^{i + 1/2} \quad \text{or} \quad \hat{v} = \gamma^i
\end{equation}
The maximum absolute error ratio between the true value $v$ and the approximation $\hat{v}$ is bounded by:
\begin{equation}
\frac{|v - \hat{v}|}{v} \le \frac{\gamma^{i+1} - \gamma^i}{\gamma^i} = \gamma - 1 = \frac{2\epsilon}{1 - \epsilon} \approx 2\epsilon
\end{equation}
By scaling the index mapping, the relative error is guaranteed to satisfy:
\begin{equation}
\text{Relative Error} \le \epsilon
\end{equation}
This bounds the error of any computed percentile (such as p50, p95, or p99) to at most 1\% of its true value.

---

\section{Tokens-Per-Output-Token (TPOT) and Token Throughput}

For streaming inference requests, latency metrics include the Time-to-First-Token ($L_{\text{ttft}}$) and the Total Latency ($L_{\text{total}}$).
Let $C$ be the number of completion tokens generated. The Tokens-Per-Output-Token ($T_{\text{tpot}}$) metric represents the average generation time per completion token:
\begin{equation}
T_{\text{tpot}} = \frac{L_{\text{total}} - L_{\text{ttft}}}{C}
\end{equation}
This metric is computed only under the constraints:
\begin{equation}
L_{\text{ttft}} > 0, \quad C > 0, \quad \text{finish\_reason} \ne \text{"timeout"}
\end{equation}
TPOT calculations are stored in a rolling Redis list of size $K=1000$ using `LPUSH` and `LTRIM` operations, presenting a space complexity of $O(K)$.

---

\section{Computational and Storage Complexity}

We analyze the update complexity for incoming spans:
\begin{enumerate}
\item **Time Complexity**: Inserting a latency value into a DDSketch is $O(\log(1/\epsilon))$ because it maps the float value to an integer bucket index and increments the bucket counter. For $\epsilon = 0.01$, this requires fewer than 500 bucket operations, completing in sub-microsecond time.
\item **Space Complexity**: The number of buckets created is bounded by the ratio of maximum to minimum values. The memory usage is:
  \begin{equation}
  \text{Memory} = O\left(\log_{\gamma}\left(\frac{v_{\text{max}}}{v_{\text{min}}}\right)\right)
  \right.
  \end{equation}
  This is constant ($O(1)$) with respect to the total number of spans $N$ processed, preventing memory leaks under high volumes.
\end{enumerate}

\end{document}
